\chapter{绪论}

\section{榜题要求与本系统的定位}

本系统面向挑战杯“揭榜挂帅”擂台赛榜题 XH-202630\cite{xh202630}。
榜题要求构建一套面向垂直领域技能培训的领域知识个性化生成与多智能体协同决策系统，
核心指标有三项：专业知识谬误率低于百分之五、
学习者画像与资源难度适配准确率不低于百分之八十五、核心知识点覆盖率不低于百分之九十。

榜题在评分标准中明确点名了两项技术要求。其一是引入“辩论与交叉验证”机制
以消除大模型在专业领域的幻觉；其二是探索大模型驱动下的动态追问与启发式交互导学，
打破静态资源的单向输入局限。这两项在本系统中分别对应第四章的双专家交叉验证
与第二章的自适应测评。

本系统选取\emph{工业机器人现场操作与调试}作为示范领域。
选择这一领域有其技术上的理由：它拥有大量\emph{封闭事实集}，
例如报警代码表、指令参数表、公差要求与保养周期。
封闭事实集的特点是对错一查便知，因而幻觉率的真值容易建立；
若选取开放性题材，幻觉率将无从测量。

\section{总体架构}

系统由五类智能体与一个编排层构成，各自职责见表 \ref{tab:agents}。
表中最后一列尤其值得注意：承担\emph{判断}职责的环节全部不使用大模型。

\begin{table}[htbp]\centering\small
\caption{各智能体的职责划分}\label{tab:agents}
\begin{tabular}{@{}p{2.2cm}p{6.4cm}p{2.4cm}@{}}
\toprule
智能体 & 职责 & 是否用大模型 \\
\midrule
学情访谈 & 自述文本转结构化背景 & 用，有规则兜底 \\
学情诊断 & 作答记录转掌握概率，盲区排序 & 不用 \\
领域专家甲 & 窄检索起草断言 & 用 \\
领域专家乙 & 宽检索起草断言 & 用 \\
交叉验证裁判 & 对齐两份断言，仲裁冲突 & 不用 \\
审核裁判 & 逐条断言核验 & 规则为主，模型为辅 \\
命题 & 现场出变式题，过五关审核 & 用 \\
决策调度 & 按正确率决定降维、巩固或进阶 & 不用 \\
\bottomrule
\end{tabular}
\end{table}

全流程的状态流转如下：

\begin{code}
INIT -> DIAGNOSE -> PLAN -> [ GENERATE -> DEBATE -> AUDIT
                              -> (REVISE -> AUDIT)* -> ASSEMBLE ]* -> READY

READY -> FEEDBACK -> DECIDE -> GENERATE ...
\end{code}

其中 \cd{GENERATE} 一步由两位领域专家并行执行，专家甲走窄检索、专家乙走宽检索；
\cd{DEBATE} 对二者产出的断言做交叉验证与仲裁；方括号表示该段对学习路径上的
每个知识点重复一次；星号表示该段可能重复零到两轮。
反馈回路从 \cd{READY} 出发，经决策调度后重新进入生成，构成动态迭代。

编排层采用\emph{手写状态机}，未采用 LangGraph 或 AutoGen 一类框架，理由有三。
其一，每条状态转移边都能直接写成单元测试，而榜题明确要求测试多智能体协同调度逻辑。
其二，出现问题时调用栈很浅，学生自己即可调试，不必先读框架源码。
其三，状态与事件由本系统自行定义，前端可视化与评测中间数据可以复用同一份结构。
待真正需要并行调度或人在回路时再行迁移，接口已经预留。

\section{三条不能破的原则}

以下三条原则是整个方案可信度的来源。后续增加功能时守住它们，
比多加两个智能体有用得多。

\subsection{大模型只管生成内容，不管控制流}

掌握度是\emph{算}出来的，教学决策是\emph{规则}判的。
若将其交给模型，同一份前测运行两遍结果不一致，
所谓“个性化适配”与“动态迭代”就变成了看运气，评委随手试两次便会露馅。

落到代码上：贝叶斯知识追踪的掌握度估计、自适应选题、追问触发条件、
降维与进阶的阈值判定，全部是算术，没有一行调用模型。
模型只负责把算好的数字翻译成自然语言，以及产出内容文本。

\subsection{没有引用的断言等于没说}

生成智能体的最小输出单元是断言，每条必须携带知识库切片编号。
审核不通过就打回重写，最多两轮，仍不通过则直接丢弃。

\emph{幻觉率能够压住，是因为架构上根本出不去，而不是生成完再去统计有多少条是错的。}
这是技术创新性评分项的核心论据。

\subsection{检索按知识点硬过滤}

生成某个知识点的资源时，模型只看得到该知识点的切片，
跨知识点的内容串味在检索阶段即被切断。
这是压制幻觉的第一道闸门，实现成本几乎为零。

\section{运行方式}

系统不需要安装任何第三方包，不需要联网，也不需要模型接口密钥。
解压后进入项目目录，执行下列命令即可验证全部功能。

首先运行一次完整闭环，观察多智能体调度的时间线。
两位学习者基础不同，生成的学习路径也完全不同：

\begin{code}
python3 cli.py P-A
python3 cli.py P-C
\end{code}

其次运行全部单元测试，当前共 237 个用例：

\begin{code}
python3 -m unittest discover -s tests -v
\end{code}

三类评测分别为三项指标的批量评测、按幻觉类型统计检出率的红队测试、
以及题库质量报告：

\begin{code}
python3 -m evalkit.run_eval --n 50
python3 -m evalkit.redteam
python3 -m evalkit.itemreport
\end{code}

消融对照是评测中最有说服力的部分，分别关闭审核环节与辩论环节，
观察各自对幻觉率的贡献：

\begin{code}
AGENTEDU_INJECT=0.3 python3 -m evalkit.run_eval --n 50 --no-audit
AGENTEDU_INJECT=0.3 python3 -m evalkit.run_eval --n 50 --no-debate
\end{code}

最后启动演示服务，浏览器访问 \cd{http://127.0.0.1:8000}：

\begin{code}
python3 server.py
\end{code}

\begin{keynote}{最重要的一个文件：web/showcase.html}
双击即可打开，不需要后端、不需要网络、不需要安装任何东西。
评委可以现场输入自己的背景、亲手答题，系统当场算出学情并生成资源。

这是专为答辩现场准备的。演示服务在别人机器上起不来是很常见的事，
端口占用、防火墙策略、Python 版本差异都可能导致失败，单文件版把这个风险彻底消掉。
后端在线时它会自动切换为实时模式，页面顶部的指示灯显示当前数据源。
\end{keynote}

\bibliographystyle{gbt7714-2005}
\bibliography{refs}
